\section{Implementation}
\label{sec:implementation}

The Magnetar reference implementation is the Go canonical at
\texttt{github.com/luxfi/magnetar} \cite{magnetargo}, tag
\texttt{v0.5.2} (commit \texttt{1c84519}). The artefact ships both
the per-validator standalone primary primitive
(\texttt{ref/go/pkg/magnetar/standalone.go} and \texttt{wire.go})
and the threshold forms documented in the appendices
(\texttt{ref/go/pkg/thbs/}, \texttt{ref/go/pkg/magnetar/combine.go}).
This section maps the per-validator standalone construction of
\S\ref{sec:construction} to the source layout.

\subsection{File layout (standalone primary)}
\label{sec:impl:layout}

\begin{table}[H]
\centering
\caption{Magnetar v0.5.2 per-validator standalone primary path
file layout. The threshold forms (\texttt{thbs/} subpackage and
\texttt{combine.go}, \texttt{dkg.go}, \texttt{threshold.go}) are
documented in Appendix~\ref{app:thbs} and Appendix~\ref{app:reveal}.}
\label{tab:magnetar-files-standalone}
\begin{tabular}{lp{8cm}}
\toprule
File & Role \\
\midrule
\multicolumn{2}{l}{\emph{Standalone primary path}} \\
\midrule
\texttt{doc.go}            & Package doc + invariant summary \\
\texttt{params.go}         & FIPS 205 parameter table; Mode dispatch \\
\texttt{types.go}          & Wire types ($\mathsf{NodeID}$, \texttt{Signature}, \texttt{PublicKey}, \texttt{PrivateKey}, \texttt{AbortEvidence}) \\
\texttt{keygen.go}         & Single-party FIPS 205 KeyGen / DeriveKey passthrough \\
\texttt{sign.go}           & Single-party FIPS 205 SignDeterministic / SignRandomized dispatch \\
\texttt{verify.go}         & FIPS 205 Verify passthrough (no Magnetar envelope) \\
\texttt{standalone.go}     & \texttt{PerValidatorKeypair}, \texttt{ValidatorSign}, \texttt{ValidatorBatchVerify}, \texttt{ValidatorAggregateCert}, \texttt{BuildAggregateCert}, \texttt{VerifyAggregateCert} \\
\texttt{aggregate.go}      & SDK-friendly \texttt{SignedBundle} / \texttt{AggregatedSignature} shape (byte-equivalent to \texttt{ValidatorAggregateCert}) \\
\texttt{wire.go}           & Canonical MAGS / MAGG envelope codec; stateless \texttt{VerifyBytes} dispatcher \\
\texttt{transcript.go}     & cSHAKE256 helpers; constant-time byte equality \\
\texttt{zeroize.go}        & Constant-time wipes for bytes, uint16 slices, PrivateKeys \\
\midrule
\multicolumn{2}{l}{\emph{Standalone primary path tests}} \\
\midrule
\texttt{standalone\_test.go}      & 7 tests: round-trip, determinism, independence, batch verify, bad-sig-counted-out, unknown-validator rejection, duplicate-dedupe \\
\texttt{aggregate\_test.go}       & Bundle / aggregate cross-checks \\
\texttt{wire\_test.go}            & MAGS / MAGG roundtrip + \textbf{TestMagnetar\_Wire\_FIPS205Verifiable} (load-bearing byte-identity gate) \\
\texttt{verify.go} / \texttt{sign\_test.go} / \texttt{keygen\_test.go} & Single-party FIPS 205 dispatch tests \\
\texttt{race\_off\_test.go} / \texttt{race\_on\_test.go} & \texttt{raceEnabled} build-tag pattern (FIPS 205 heavy paths skipped under race) \\
\midrule
\multicolumn{2}{l}{\emph{KAT generator}} \\
\midrule
\texttt{cmd/genkat/main.go}     & Deterministic KAT regeneration binary \\
\bottomrule
\end{tabular}
\end{table}

The strict layering invariant: every standalone-path file's
package-public surface calls only into lower-level packages (FIPS
205 via \texttt{circl}, cSHAKE256); no cross-file circular
dependency. The standalone primitive is implemented top-down with
\texttt{VerifyAggregateCert} sitting at the highest layer, calling
\texttt{ValidatorBatchVerify}, which dispatches the
\texttt{circl.slhdsa.Verify} workers via the same goroutine
fork-join pattern as \texttt{luxfi/crypto/slhdsa.VerifyBatch}.

\subsection{Dependencies}
\label{sec:impl:deps}

The dependency graph is deliberately narrow.

\begin{itemize}[leftmargin=2em,nosep]
  \item \texttt{github.com/cloudflare/circl}~v1.6.3 \cite{circl}:
        FIPS 205 SLH-DSA reference (BSD-3-Clause). Magnetar uses
        \texttt{circl/sign/slhdsa} for
        $\textsc{KeyGen}$ / $\textsc{DeriveKey}$ /
        $\textsc{SignDeterministic}$ / $\textsc{SignRandomized}$ /
        $\textsc{Verify}$. The \texttt{circl} implementation is the
        most widely-deployed Go FIPS 205 backend; its constant-time
        posture has been audited under Cloudflare's internal review
        and is used in \texttt{cloudflare/cf-mldsa-precompile} in
        production.
  \item \texttt{golang.org/x/crypto}~v0.32.0: provides
        \texttt{sha3.NewCShake256} for cSHAKE256 and the SP 800-185
        encoding primitives. Standard-library-adjacent dependency.
  \item \texttt{golang.org/x/sys}~v0.29.0 (indirect): platform
        syscall layer used by \texttt{crypto/rand}.
\end{itemize}

The dependency on \texttt{circl} is the load-bearing one. The
byte-identity claim of
$\textsc{TestMagnetar\_Wire\_FIPS205Verifiable}$ binds Magnetar's
output to the exact \texttt{circl}-emitted bytes. If \texttt{circl}
v1.6.3 were ever deprecated or its FIPS 205 implementation found
non-conformant, Magnetar would need to switch to a different FIPS 205
backend. Two alternative backends are candidate: a future Go
\texttt{crypto/slhdsa} once Go's \texttt{crypto} package gains FIPS
205 support, and a vendored Magnetar copy of the NIST FIPS 205
reference C implementation wrapped via cgo. Both are tracked in
\texttt{TRUSTED-COMPUTING-BASE.md}.

\subsection{Build and test discipline}
\label{sec:impl:build}

The reproducible build is:
\begin{lstlisting}[language=bash,caption={Magnetar reproducible build.}]
cd ref/go
GOWORK=off go build ./...
GOWORK=off go test -short ./pkg/magnetar/...
GOWORK=off go test -count=1 -run TestMagnetar_Wire_FIPS205Verifiable ./pkg/magnetar/...
\end{lstlisting}

The \texttt{GOWORK=off} flag is required to suppress a workspace
file that overlays unrelated Lux packages; without it, Go's module
loader pulls in transitive paths that may not be on the developer's
machine. (This is a long-standing Lux-internal annoyance, not a
Magnetar-specific quirk.)

The \texttt{-count=1} flag bypasses Go's test cache and forces a
fresh execution, which is the standard discipline for crypto
regression gates: a cached PASS on a stale binary is not evidence.

The \texttt{-short} flag selects the fast CI configuration; the
SLH-DSA-heavy test paths run only the cheap configurations under
\texttt{-short}.

\paragraph{Race detector pattern.} The \texttt{raceEnabled}
build-tag pattern at \texttt{race\_off\_test.go} /
\texttt{race\_on\_test.go} gates SLH-DSA-heavy tests behind
\texttt{!race}; the race detector adds 5--10x overhead on the
SLH-DSA hot path, and the tests are \emph{single-goroutine} so they
cannot detect a race regardless. The
$\textsc{TestMagnetar\_Wire\_FIPS205Verifiable}$ gate uses
\texttt{skipUnderRace} to skip under \texttt{-race}; race-on runs
cover the cheap wire-codec roundtrip tests. This pattern is shared
with Pulsar and Corona and is documented in their respective
READMEs.

\paragraph{KAT regeneration.} The KAT generator at
\texttt{cmd/genkat} accepts a single 48-byte master seed
(\texttt{masterSeedHex}) and produces deterministic vector files.
Re-running the binary on a clean checkout produces byte-identical
output --- this is the deterministic-fixture gate. Drift between
committed JSON and freshly-regenerated output is a CI failure.

\subsection{Headline regression: TestMagnetar\_Wire\_FIPS205Verifiable}
\label{sec:impl:wirefips}

The byte-identity claim is gated by the test at
\texttt{ref/go/pkg/magnetar/wire\_test.go:124-220}. The test routes
through the v0.5 per-validator standalone primary path:

\begin{lstlisting}[language=go,caption={Excerpt: \texttt{TestMagnetar\_Wire\_FIPS205Verifiable} core assertions.}]
sk, pk, _ := PerValidatorKeypair(params, newDetReader(seed))
sigBytes, _ := ValidatorSign(sk, nil, msg)
sig := &Signature{Mode: mode, Bytes: sigBytes}

// Frame both via the canonical wire codec.
gkWire, _  := MarshalGroupKey(pk)
sigWire, _ := sig.MarshalBinary()

// Strip the 11-byte header to recover the raw FIPS 205 payload.
fipsPK  := extractFIPS205Payload(t, gkWire, wireMagicMagnetarGroupKey)
fipsSig := extractFIPS205Payload(t, sigWire, wireMagicMagnetarSig)

// Verify under cloudflare/circl directly. No magnetar code path
// is hit on the verify side. This is the load-bearing assertion.
if !verifyDirectCircl(mode, fipsPK, msg, fipsSig) {
  t.Fatalf("FIPS 205 Verify rejected the unwrapped Magnetar bytes")
}
\end{lstlisting}

The test runs for all three SHAKE parameter sets (SHAKE-192s,
SHAKE-192f, SHAKE-256s) and passes under \texttt{-count=1} and
\texttt{-race}. It is the single test that pins
\S\ref{sec:security:byte-equality} empirically.

\subsection{Configuration-as-data}
\label{sec:impl:config}

The Mode-to-Params dispatch in \texttt{params.go} is a static table
with no runtime logic:
\begin{lstlisting}[language=go,caption={Mode dispatch in \texttt{params.go}.}]
var ParamsM192s = &Params{
    Mode:           ModeM192s,
    N:              24,
    PublicKeySize:  48,
    PrivateKeySize: 96,
    SignatureSize:  16224,
    SeedSize:       96,
    ID:             slhdsa.SHAKE_192s,
    ShamirPrime:    257,
}
\end{lstlisting}

The \texttt{Params.Validate()} method checks that a supplied
\texttt{*Params} matches one of the canonical tables exactly. Any
deviation --- a tampered \texttt{ShamirPrime}, a wrong
\texttt{SeedSize}, an unsupported \texttt{Mode} --- is rejected at
the entry point of \texttt{PerValidatorKeypair},
\texttt{ValidatorSign}, \texttt{ValidatorBatchVerify}, and
\texttt{VerifyAggregateCert}. The \texttt{*Params} argument is the
single point where parameter selection enters the crypto path; this
matches the ``one and only one way to do everything'' discipline.

\subsection{Concurrency posture}
\label{sec:impl:concurrency}

The standalone primary path's
\texttt{ValidatorBatchVerify} dispatches across
$\min(\mathit{GOMAXPROCS}, N)$ goroutines above the
\texttt{verifyAggregatedConcurrentThreshold} constant (currently
$\sim 16$); below that, the batch runs serial. Each worker
invokes \texttt{circl.slhdsa.Verify} on a disjoint range of the
input slices; there is no shared mutable state.

\texttt{ValidatorSign} is strictly single-goroutine: the FIPS 205
signing path is sequential per signature, and a deployment running
many concurrent signing sessions allocates one signer goroutine per
session. The threshold forms documented in Appendix~\ref{app:thbs}
and Appendix~\ref{app:reveal} are also strictly single-goroutine
per session.

\subsection{Lux Polaris integration}
\label{sec:impl:polaris}

Magnetar plugs into the Lux Polaris cert profile at the LSS
\cite{luxlssmpc} adapter layer. The integration is intentionally
mirror-symmetric with Pulsar and Corona at the public-BFT path: a
Polaris signer is the composition $\textsc{PolarisSigner} =
(\textsc{PulsarSigner}, \textsc{CoronaSigner},
\textsc{MagnetarSigner})$, with the three lanes signing the same
Prism-bound transcript. The Magnetar lane's
$\textsc{MagnetarSigner}$ is the per-validator standalone primitive
of this paper.

The activation cert composition: a Polaris generation transition
requires three independent per-lane operations in lockstep.
For Magnetar, this is per-validator rotation of $pk_i$ in the
registry; the standalone primitive does not have a DKG-style
``activation cert'' construction (it does not need one --- no
shared key to activate). The LSS adapter contract
\cite{luxlssmpc} is primitive-agnostic at its boundary; Magnetar's
standalone adapter implements an empty $\textsc{Reshare}$
(rotation is per-validator, not committee-wide) and a
$\textsc{KeyRotate}$ that signs the old $pk_i$ replacing under
$sk_i^{\mathrm{old}}$ before installing $(sk_i^{\mathrm{new}},
pk_i^{\mathrm{new}})$.

For the threshold forms (Appendix~\ref{app:thbs},
Appendix~\ref{app:reveal}), the LSS lifecycle is committee-wide
and matches the Pulsar / Corona threshold adapters. Those are
appropriate for the M-Chain bridge custody and A-Chain
confidential compute paths.
